\section{Resumen}
\begin{abstract}
En esta práctica se abordó el modelado geométrico, es decir, la construcción de
objetos tridimensionales a partir de primitivas más sencillas. Se
implementaron con OpenGL 3.3 las tres primeras primitivas que solicita el
manual: un plano paralelo al plano $XY$ de $2\times2$ unidades, un cubo de
$2\times2\times2$ unidades y un círculo de radio unitario, las tres centradas
en el origen del sistema de referencia y descritas mediante listas de vértices
e índices que la clase \texttt{Mesh} envía a la unidad de procesamiento
gráfico. Cada figura puede observarse desde dos puntos de vista distintos,
asignados a las teclas \texttt{F1} y \texttt{F2}, y puede alternarse entre
relleno y alambre para exhibir la triangulación con la que fue construida.
El trabajo dejó ver que una primitiva geométrica no se dibuja directamente:
antes se descompone en primitivas gráficas, en este caso triángulos, y de la
forma en que se elija esa descomposición dependen tanto el número de vértices
como la apariencia final. Al ejecutar la clase \texttt{Circle} del proyecto
base detectamos que dejaba sin inicializar el color de casi todos sus vértices y que su
arreglo de índices rebasaba el tamaño del arreglo de vértices, por lo que se
reescribió como un abanico de triángulos con un único vértice central.
\end{abstract}

\renewcommand{\abstractname}{Abstract}
\begin{abstract}
This laboratory session deals with geometric modelling, that is, the
construction of three-dimensional objects out of simpler primitives. The first
three primitives requested by the manual were implemented in OpenGL 3.3: a
plane parallel to the $XY$ plane measuring $2\times2$ units, a cube measuring
$2\times2\times2$ units and a circle of unit radius, all three centred at the
origin of the reference system and described by vertex and index lists that
the \texttt{Mesh} class uploads to the graphics processing unit. Each figure
can be observed from two different viewpoints, bound to the \texttt{F1} and
\texttt{F2} keys, and can be switched between filled and wireframe rendering
in order to expose the triangulation it was built from. The work made clear
that a geometric primitive is not drawn directly: it is first decomposed into
graphic primitives, triangles in this case, and both the vertex count and the
final appearance depend on how that decomposition is chosen. When running the
\texttt{Circle} class supplied with the base project we found that it left the colour of nearly all its vertices uninitialised and
that its index array exceeded the size of the vertex array, so it was rewritten as a triangle fan with a single
central vertex.
\end{abstract}
